\sezione{10}{Esempio 3: regressione}{Parte III · RandomForestRegressor e confronto onesto}

\begin{inbreve}
Ora il target è un \hl{numero}: una misura della progressione del diabete a un anno, per 442 pazienti descritti da 10 variabili (età, sesso, BMI, pressione e sei analisi del sangue; dataset \emph{Diabetes} di scikit-learn). Confrontiamo il Random Forest con una regressione lineare e con un albero singolo, e verifichiamo che la previsione della foresta è \key{la media delle previsioni degli alberi}.
\end{inbreve}

\subsection{Il codice}

\codicefile{es3_regressione.py}{1}{44}

\subsection{Spiegazione}

\begin{tecnica}[Il dizionario di modelli \ \riga{17--23}]
Mettere i modelli in un dizionario \texttt{\{nome: modello\}} permette di addestrarli e valutarli con lo \key{stesso ciclo} (\riga{24--28}): stesso split, stesse metriche, confronto equo. Per il Random Forest si usa la regola classica della regressione, $m=p/3$ (\texttt{max\_features=1/3}), e foglie di almeno 5 esempi (\texttt{min\_samples\_leaf=5}), come suggerito in letteratura.
\end{tecnica}

\begin{tecnica}[Le metriche di regressione \ \riga{27--28}]
\[
\mathrm{MAE}=\frac1n\sum_i|y_i-\hat y_i|,\qquad
\mathrm{RMSE}=\sqrt{\frac1n\sum_i(y_i-\hat y_i)^2},\qquad
R^2=1-\frac{\sum_i(y_i-\hat y_i)^2}{\sum_i(y_i-\bar y)^2}.
\]
MAE e RMSE sono nell'unità del target (più bassi è meglio). $R^2$ confronta il modello con chi prevede sempre la media: $R^2=0$ vale quanto la media, $R^2<0$ è \key{peggio} della media.
\end{tecnica}

\begin{tecnica}[La previsione è una media \ \riga{34--37}]
\texttt{rf.estimators\_} contiene i 500 alberi. Si chiede a ciascuno la previsione per lo stesso paziente e se ne fa la media: il risultato coincide con \texttt{rf.predict}. Gli alberi interni sono stati addestrati su array NumPy senza nomi di colonna, per questo si passa \texttt{x0.values}.
\end{tecnica}

\begin{tecnica}[Confronto robusto \ \riga{40--44}]
Un solo split train/test può essere fortunato o sfortunato. \texttt{RepeatedKFold(n\_splits=5, n\_repeats=3)} ripete la validazione a 5 fold 3 volte con mescolamenti diversi: 15 valutazioni per modello, di cui si riportano media e deviazione standard.
\end{tecnica}

\begin{terminale}{python es3_regressione.py}
Regressione lineare  MAE= 45.12  RMSE= 56.39  R2=0.359
Albero singolo       MAE= 60.42  RMSE= 76.59  R2=-0.182
Random Forest        MAE= 46.40  RMSE= 58.89  R2=0.301
R2 out-of-bag: 0.508
media alberi = 234.902   rf.predict = 234.902
dispersione tra alberi (dev. std) = 52.76
Regressione lineare  R2 CV = ¤0.490 +/- 0.066¤
Albero singolo       R2 CV = ¤-0.170 +/- 0.221¤
Random Forest        R2 CV = ¤0.465 +/- 0.085¤
\end{terminale}

\begin{esempio}[Interpretazione: una lezione importante]
\begin{itemize}
  \item \hl{L'albero singolo è peggio della media} ($R^2<0$): un albero profondo su 331 esempi rumorosi fa overfitting pesante. La foresta, con gli stessi alberi messi in media, passa da $-0.17$ a $+0.47$ in validazione incrociata: è la \key{riduzione di varianza} del capitolo 4 vista dal vivo.
  \item \hl{La regressione lineare vince di poco} (0.490 contro 0.465, differenza minore della deviazione standard). Su questo dataset la relazione è quasi lineare e i dati sono pochi: il Random Forest non ha nulla da guadagnare dalla sua flessibilità. \key{Non è sempre il modello migliore}: confrontalo sempre con un modello semplice.
  \item \hl{OOB 0.508 contro test 0.301}: con soli 111 pazienti nel test la stima di un singolo split è molto variabile. La CV ripetuta (0.465) sta nel mezzo: per dataset piccoli fidati di più della CV.
  \item La media delle previsioni dei 500 alberi è esattamente \texttt{rf.predict} (234.902), ma i singoli alberi sono molto dispersi (deviazione standard 52.8).
\end{itemize}
\end{esempio}

\par\noindent\begin{minipage}{\linewidth}
\fig[.8\linewidth]{dispersione.pdf}
\didascalia{Previsioni dei 500 alberi per il primo paziente del test. La foresta restituisce la media (234.9); il valore vero è 321. Ogni albero è molto impreciso, ma la media è stabile. La dispersione dà anche un'idea (approssimativa) dell'incertezza.}
\end{minipage}\par

\begin{attenzione}[Il Random Forest non estrapola]
In regressione ogni foglia prevede la \hl{media dei target} degli esempi che contiene. Quindi la foresta non può mai prevedere un valore fuori dall'intervallo dei target visti in addestramento: se il fenomeno ha un trend (prezzi che crescono negli anni, serie temporali), il Random Forest «si appiattisce» oltre i dati noti. In questi casi va affiancato o sostituito da modelli che estrapolano.
\end{attenzione}

\begin{appunti}{Riepilogo del capitolo}
\begin{itemize}
  \item \texttt{RandomForestRegressor}: stessa API; criterio \texttt{squared\_error}; previsione = media degli alberi.
  \item In sklearn il default è \texttt{max\_features=1.0}: per la regola classica scrivi \texttt{1/3}.
  \item Metriche: MAE, RMSE, $R^2$. Confronta sempre con un modello semplice e con CV ripetuta.
  \item La foresta riduce enormemente la varianza dell'albero singolo ma non estrapola fuori dal range dei dati.
\end{itemize}
\end{appunti}
